% PhD application version adapted from the MLE resume.
% Compile: tectonic Yuanbo_Pang_CV_PhD_2027.tex
\documentclass[9pt, letterpaper]{extarticle}

\usepackage[top=0.4in,bottom=0.4in,left=0.55in,right=0.55in]{geometry}
\usepackage{enumitem}
\usepackage{titlesec}
\usepackage{hyperref}

\pagestyle{empty}
\setlength{\parindent}{0pt}
\setlength{\parskip}{0pt}

\hypersetup{
  colorlinks=true,
  urlcolor=black,
  pdfauthor={Yuanbo Pang},
  pdftitle={Yuanbo Pang CV}
}

\titleformat{\section}
  {\large\scshape\bfseries\raggedright}
  {}{0em}{}[\vspace{-2pt}\titlerule]
\titlespacing*{\section}{0pt}{7pt}{4pt}

\newcommand{\role}[2]{%
  \vspace{3pt}%
  \noindent\textbf{#1}\hfill\textit{#2}\par\vspace{1pt}%
}
\newcommand{\rolesub}[1]{%
  \noindent\textit{#1}\par\vspace{1pt}%
}

\setlist[itemize]{
  leftmargin=1.2em,
  itemsep=1.6pt,
  parsep=0pt,
  topsep=1.6pt,
  partopsep=0pt
}

\begin{document}

\begin{center}
  {\LARGE\bfseries YUANBO PANG}\\[2pt]
  {\large UC Berkeley EECS \quad$\vert$\quad Foundation Agents \& Robotics}\\[2pt]
  \href{mailto:yuanbopang@berkeley.edu}{yuanbopang@berkeley.edu}
  $\;\cdot\;$ (+1) 408-460-7526
  $\;\cdot\;$ \href{https://www.linkedin.com/in/yuanbopang/}{LinkedIn}
  $\;\cdot\;$ \href{https://github.com/yuanboP}{GitHub}
\end{center}

\section{Education}
\role{University of California, Berkeley}{Expected Dec 2026}
\rolesub{B.S.\ Electrical Engineering \& Computer Sciences \quad GPA: 3.762/4.0 (in-major: 3.811)}
\rolesub{Relevant coursework: CS 170, CS 168, CS 184, EECS 106A, CS 61B, CS 70}
\role{Stanford University (Summer Session) --- CS 229: Machine Learning}{Summer 2025}
\role{De Anza College --- Computer Science, GPA: 4.0/4.0}{2023 -- 2025}

\section{Research Experience}

\role{Research Assistant, Berkeley RDI / BAIR (Prof.\ Dawn Song)}{Oct 2025 -- Present}
\rolesub{Agents' Last Exam --- co-first author on V1; leading two research tracks in V2}
\begin{itemize}
  \item \textbf{Foundation Agent / 3D spatial understanding.} Building training and evaluation environments in simulation (Blender, NVIDIA Isaac Lab) where agents actively move the viewpoint to explore a scene and judge physical state --- 3D understanding through exploration, not inference from a fixed 2D view.
  \item \textbf{Ultra-long-horizon engineering.} Constructing tasks from large-scale production software that human teams measure in months to years, with a harness that evaluates \emph{teams} of agents --- coordination, division of labor, and context handoff at horizons far beyond single-agent benchmarks.
  \item For V1, built task/environment pipelines, computer-use and tool interfaces, scalable GCP execution, and programmatic evaluators producing verifiable rewards for RLVR; now run ALE evaluations with frontier model teams (OpenAI, Anthropic, Meta TBD Lab, Qwen, Seed), analyzing agent trajectories and reporting failure modes back to the labs.
\end{itemize}

\role{Robotics Research, Prof.\ Masayoshi Tomizuka's Lab, UC Berkeley}{2026 -- Present}
\rolesub{Large-Scale Robot-Learning Data Benchmark --- mentored by Zhuo Xu}
\begin{itemize}
  \item Built \texttt{roboqa}: dataset-quality benchmark at EgoDex scale (338K episodes, 829 h) --- health, diversity, and trainability scored separately; every score published with full provenance.
  \item Agent workflows inside processing and annotation: VLM-judged quality audits; raw captures auto-converted to a canonical, embodiment-tagged format (egocentric hand, gripper, humanoid).
  \item Train and post-train VLAs on 500\,h of curated data; score the resulting policies in simulation.
\end{itemize}

\role{URAP Research Assistant, Prof.\ Avideh Zakhor, UC Berkeley EECS}{Spring 2026 -- Present}
\begin{itemize}
  \item Building a high-quality UAV-to-UAV dataset and PyTorch pipelines for zero-shot detection and multi-object tracking under high-speed ego-motion; integrating optical flow with YOLO and SAM2-style baselines. Targeting CVPR submission.
\end{itemize}

\section{Engineering Experience}

\role{Founder, Equile --- ``The dry lab for AI scientists''}{Apr 2026 -- Oct 2026}
\begin{itemize}
  \item Sole owner of the runtime and control plane: serverless hosting for agent harnesses (Claude Code, Codex, Cursor, Grok) --- isolated per-session sandboxes, multi-tenant packing that cuts per-agent compute $>$20$\times$, snapshot pause/resume, public SDK/CLI.
\end{itemize}

\role{Agentic Tech Staff, Starbot}{Nov 2025 -- May 2026}
\begin{itemize}
  \item Led the agentic stack for autonomous service robots: vision-language-action models, agentic planning and control, multi-robot coordination.
  \item Shipped as the CES 2026 demo system, deployed on the show floor 8 hours/day across coffee preparation, popcorn making, and front-desk reception; covered live by Japan's NHK News, demoed to 1{,}500+ attendees, and converted into 50+ inbound customer conversations.
\end{itemize}

\role{Software Engineering Intern, Geopogo}{Fall 2025}
\begin{itemize}
  \item Shipped a production multimodal-generation workflow (React/TypeScript, FastAPI, Gemini, Firebase, CI/CD).
\end{itemize}

\section{Research Interests}
\textbf{Foundation Agents:} 3D and environment understanding; RL systems for agent--environment interaction\\[2pt]
\textbf{Robotics:} 3D understanding to strengthen world-action models (VLA)

\section{Publications}
% TODO: swap "Under review at" to "To appear at" the day the acceptance notification lands.
\textbf{Agents' Last Exam.}\\[1pt]
Yiyou Sun\textsuperscript{*}, Xinyang Han\textsuperscript{*}, Weichen Zhang\textsuperscript{*}, \textbf{Yuanbo Pang}\textsuperscript{*}, et al., Dawn Song. \textit{Under review at NeurIPS 2026.} \href{https://arxiv.org/abs/2606.05405}{arXiv:2606.05405}. \quad(\textsuperscript{*}Co-first author \& core contributor; 310 authors total)\\[1pt]
\textit{\#1 Paper of the Week, Hugging Face Daily Papers (Jun 2026).}\\[6pt]
\textbf{ScienceIDE: Turning World's Scientific Codebase into Agent Learnable Environments.}\\[1pt]
Hejia Geng\textsuperscript{*}, Zesen Huang\textsuperscript{*}, Haoyang Li, Wenbin Li, Koutian Wu, Zihan Zhou, \textbf{Yuanbo Pang}, et al., Ling Yang. \href{https://arxiv.org/abs/2609.19134}{arXiv:2609.19134}. \quad(AItonomy Foundation; sponsors: PhAI-Labs, Qwen)\\[1pt]
\textit{\#2 Paper of the Day, Hugging Face Daily Papers (Sep 2026).}\\[6pt]
\textbf{Agent's GameDev Exam.}\\[1pt]
Zonglin Di, Yinhe Wang, \textbf{Yuanbo Pang}, Yiyou Sun, Weichen Zhang, Xinyang Han, et al., Dawn Song. \textit{Under double-blind review at ICLR 2027.}

\section{Awards}
HackMIT Modal Sponsor Prize (2025); Smartest AI Agent, HackMIT (2024).

\section{Technical Skills}
\textbf{ML / RL:} PyTorch, reinforcement learning (RLVR), VLA, model finetuning \& evaluation\\[0.5pt]
\textbf{Vision / Robotics / Sim:} YOLO, SAM2, optical flow, Isaac Lab, Blender\\[0.5pt]
\textbf{Agent / Systems:} benchmarks, automated evaluators, tool-use \& MCP, Docker, GCP, Modal, FastAPI, Linux

\end{document}
